\scanpage{35}
\deleted{\uline{Corollaire} Soit $G$ lisse sur $S$, et supposons que}

\deleted{On sait qu'un tel ss-gpe \unsure{de|est} $S$}

Appliquons ceci au cas de $C=G_0$; alors, si $S$ est \deleted{\unsure{loc}} str. local, les $ {}_nC$ se prolongent de \uline{façon unique} en des sous-groupes $H(n)$ \deleted{de} étales finis de $G$, \unsure{nécessairement|nécessairemt} centraux, et
\[
{}_mH(n)=H(m)\qquad\text{si }m\mid n
\]

Noter que $\rho_{ab}$ et $\rho_r$, donc $\rho_{ab}+\rho_r=\rho_{\mathrm{rig}}$ sont fonctions croissantes d'un groupe \deleted{lisse} alg.

Soit $K_1\deleted{\unsure{1}}$ l'adhérence de la \unsure{suite|réun$^{\rm ion}$|limite} des $H(n)_1$, \inserted{[qu'on peut supposer lisse$/k(s_1)$ !]} et $\unsure{K|K_1|\overline{K}_1}=\overline{K_1}^{\circ}$, alors \deleted{\unsure{Pour} \unsure{avoir} \unsure{le} \unsure{traiter}}.

Alors
\[
\begin{aligned}
\rho_{ab}(K_1)&\leq\rho_{ab}(G_1)=\rho'_{ab}\\
\rho_{\mathrm{rig}}(K_1)&\leq\rho_{\mathrm{rig}}(\deleted{\unsure{K}}G_1)=\rho'_{\mathrm{rig}}.
\end{aligned}
\]
\inserted{D'autre part, on voit facilement \inserted{grâce à $K_0\supset H(n)_0={}_nG_0$} (\unsure{lemme|comme} ....),}
\[
\begin{aligned}
\inserted{\rho_{ab}(G_0)=\rho_{ab}(K_0)}\\
\inserted{\rho_r(G_0)=\rho_r(K_0)}
\end{aligned}
\]
Donc pour prouver (1) et (2), on est ramené à prouver la $\overline{\mathrm m}$ chose pour $K$ remplaçant $G$.

On a perdu la lissité de $G$, mais on a gagné la \uline{commutativité}, qui implique d'ailleurs $\rho_n=\deleted{\unsure{\rho_n}}\rho'_n=d$. L'inégalité (2) équivaut alors à l'inégalité (6), \deleted{donc on} \inserted{et}
\margin{$\text{\textcircled{\scriptsize 1}}\ \text{à}\ \text{\textcircled{\scriptsize 7}}\ \text{\unsure{n.|ou|et}}\ \text{\textcircled{\scriptsize 8}}$}
% archive sheet number: 34
